\documentclass[11pt,a4paper]{article}
\usepackage[T1]{fontenc}
\usepackage{lmodern}
\usepackage[margin=2.5cm]{geometry}
\usepackage{amsmath}
\usepackage{graphicx}
\usepackage{booktabs}
\usepackage{authblk}
\usepackage[round]{natbib}
\usepackage[hidelinks]{hyperref}

\title{Title of the Paper}
\author[1]{First Author}
\author[1,2]{Second Author}
\affil[1]{Department of Biology, University, City, Country}
\affil[2]{Institute, City, Country}
\date{}

\begin{document}

\maketitle

\begin{abstract}
State the problem, what you did, and the main result in a few sentences.
\end{abstract}

\noindent\textbf{Keywords:} first keyword, second keyword, third keyword

\section*{Introduction}
Introduce the biological question and why it matters. Cite in parentheses
\citep{watson1953}, or by name, as \citet{altschul1990} showed.

\section*{Results}
\subsection*{First finding}
Lead each subsection with its finding (Figure~\ref{fig:overview}).

\begin{figure}[htbp]
  \centering
  \includegraphics[width=0.6\linewidth]{example-image}
  \caption{\textbf{A one-sentence title.} What each panel shows. Upload an image and use its file name here.}
  \label{fig:overview}
\end{figure}

\subsection*{Second finding}
Continue with the next result~\citep{jumper2021}.

\section*{Discussion}
Interpret the results, compare them with earlier work and state the limitations.

\section*{Methods}
\subsection*{Experimental model}
Describe organisms, cell lines and their sources.

\subsection*{Quantification and statistical analysis}
Name the tests, the number of replicates and what error bars show.

\section*{Data and code availability}
Give accession numbers and repositories.

\section*{Acknowledgements}
Thank the people and funding that supported this work.

\bibliographystyle{plainnat}
\bibliography{references}

\end{document}
